\section{Related Work}

\paragraph{Time-series classification benchmarks.}
The UCR archive is a standard benchmark suite for univariate time-series classification, with a broad collection of small-to-medium datasets used across many methodological studies \citep{dau2019ucr}. Python toolkits such as aeon provide practical infrastructure for time-series learning and archive access \citep{middlehurst2024aeon}. Recent model work continues to explore compact convolutional, transformer, and hybrid architectures for time-series classification \citep{lan2024ctctime}. Our work does not aim to introduce a new classifier. Instead, it wraps a compact classifier with a post-hoc uncertainty layer and focuses on the calibration behavior under corruptions.

\paragraph{Conformal prediction.}
Conformal prediction constructs prediction sets from nonconformity scores and calibration data \citep{vovk2005algorithmic}. Its appeal in modern machine learning is that it can wrap black-box predictors and provide distribution-free coverage under exchangeability \citep{angelopoulos2021gentle}. In classification, prediction-set size is a natural efficiency measure: a method that attains the target coverage with smaller sets is more informative. Adaptive prediction sets and regularized adaptive prediction sets change the classification score to a cumulative probability score and are strong natural baselines for efficient conformal classification \citep{romano2020classification,angelopoulos2021uncertainty}. Mondrian conformal prediction conditions calibration on a discrete taxonomy, providing a useful contrast to our soft fingerprint weighting. Our method uses the standard split-conformal classification score $1-p_\theta(y\mid x)$, but changes how the calibration threshold is estimated for corrupted test examples.

\paragraph{Weighted and local conformal methods.}
Our soft weighting is closely related to weighted conformal prediction under covariate shift and localized conformal prediction \citep{tibshirani2019covariate,guan2023localized}. Those lines of work clarify that weighting can be principled when the relevant density ratio or localization rule is well specified. \fastcp{} is more modest: the weights are heuristic similarities in a perturbation-fingerprint space, not estimated likelihood ratios, and therefore they do not inherit the coverage guarantees of covariate-shift conformal prediction. This distinction is central to our framing.

\paragraph{Corrupted and out-of-distribution time series.}
Nanopoulos and Buza study conformal prediction for out-of-distribution time-series classification with distortions such as gaps \citep{nanopoulos2025conformal}. Their work motivates the specific problem setting we consider: common production distortions can make neural classifiers overconfident or poorly calibrated. Our contribution is complementary. Rather than focusing only on global calibration variants, we compute perturbation fingerprints and use them to locally weight calibration examples.

\paragraph{Robust conformal methods.}
Franco et al. combine conformal prediction with randomized smoothing for robustness against real-world time-series perturbations \citep{franco2024guaranteeing}. This line of work is stronger when formal robustness arguments or smoothing certificates are required. \fastcp{} targets a different regime: fast post-hoc calibration under limited compute. It does not replace certified smoothing methods, but provides a cheaper empirical alternative when repeated transformed inference is not practical.

\paragraph{Adaptive inference and uncertainty.}
Early-exit networks and nested prediction sets study uncertainty across intermediate exits of a model \citep{jazbec2024early}. Such work modifies the architecture or analyzes sequential predictions within a network. \fastcp{} is orthogonal: it leaves the classifier fixed and adapts only the conformal calibration threshold based on observed perturbation fingerprints.
